\paragraph{Search mode and cost function (Table~\ref{tab:search}).} Greedy decoding of the transformer is much better than best-first search on its cumulative $-\log p$ at every size (solved at 41--70: \rhval{abl_search/transformer-greedy/test_41_70/solved/mean:2} vs.\ \rhval{main/transformer-policy/test_41_70/solved/mean:2}; paired difference \rhval{analysis/analysis/test_41_70/d_greedy_tf_t41/mean:2}, $p=$ \rhval{analysis/analysis/test_41_70/p_greedy_tf_t41/mean}; at 21--40 \rhval{analysis/analysis/test_41_70/d_greedy_tf_t21/mean:2}, $p=$ \rhval{analysis/analysis/test_41_70/p_greedy_tf_t21/mean}; at 11--20 \rhval{analysis/analysis/test_41_70/d_greedy_tf_t11/mean:2}, $p=$ \rhval{analysis/analysis/test_41_70/p_greedy_tf_t11/mean}). \textbf{H4 is therefore refuted}: backtracking search with $-\log p$ cost is worse, not better, than following the policy. Replacing only the cost by the policy's rank (same trained model) gives \rhval{abl_search/transformer-rank-cost/test_41_70/solved/mean:2} at 41--70, statistically indistinguishable from the heuristic there (paired $p=$ \rhval{analysis/analysis/test_41_70/p_tfrank_heur_t41/mean}; difference \rhval{analysis/analysis/test_41_70/d_tfrank_heur_t41/mean:3}). Because the model and training are identical in these runs and only the step cost changes, the loss between the main transformer rows and these rows is attributable to the cost function. We did not test why $-\log p$ is a poor cost; one plausible reason (not isolated here) is that cumulative negative log-probability grows with proof depth, so a long proof of $\approx$ tens of steps with small per-step costs competes with many shallow alternatives, which wastes expansions under a tight budget.

\begin{table}[t]\centering\caption{Search-mode and cost ablations (post-hoc variants of the main systems; \emph{rank cost}: step cost is the tactic's rank under the policy; \emph{greedy}: top-1 only). Same trained models as the main table.}\label{tab:search}
\resizebox{\columnwidth}{!}{\begin{tabular}{llccc}
\toprule
Method & Task & solved $\uparrow$ & mean\_exp $\downarrow$ & $n$ \\
\midrule
Transformer greedy & test\_11\_20 & \underline{1.00 $\pm$ 0.00} & 11.89 $\pm$ 0.31 & 5 \\
Transformer rank cost & test\_11\_20 & 1.00 $\pm$ 0.00 & \underline{11.39 $\pm$ 0.29} & 5 \\
MLP rank cost & test\_11\_20 & \textbf{1.00 $\pm$ 0.00} & \textbf{9.07 $\pm$ 0.30} & 5 \\
\midrule
Transformer greedy & test\_21\_40 & \underline{0.94 $\pm$ 0.02} & 27.96 $\pm$ 2.93 & 5 \\
Transformer rank cost & test\_21\_40 & 0.94 $\pm$ 0.02 & \underline{26.60 $\pm$ 3.27} & 5 \\
MLP rank cost & test\_21\_40 & \textbf{0.99 $\pm$ 0.00} & \textbf{19.51 $\pm$ 2.43} & 5 \\
\midrule
Transformer greedy & test\_41\_70 & 0.73 $\pm$ 0.04 & 45.84 $\pm$ 3.17 & 5 \\
Transformer rank cost & test\_41\_70 & \underline{0.73 $\pm$ 0.03} & \underline{44.58 $\pm$ 3.36} & 5 \\
MLP rank cost & test\_41\_70 & \textbf{0.88 $\pm$ 0.04} & \textbf{33.77 $\pm$ 3.24} & 5 \\
\bottomrule
\end{tabular}
}
\end{table}

\paragraph{Post-hoc finding: MLP with rank cost.} The feature MLP with rank cost solves \rhval{abl_search/mlp-rank-cost/test_41_70/solved/mean:2} at 41--70, \rhval{abl_search/mlp-rank-cost/test_21_40/solved/mean:2} at 21--40 and \rhval{abl_search/mlp-rank-cost/test_11_20/solved/mean:2} at 11--20, against the heuristic's \rhval{main/hand-heuristic/test_41_70/solved/mean:2}, \rhval{main/hand-heuristic/test_21_40/solved/mean:2} and \rhval{main/hand-heuristic/test_11_20/solved/mean:2}. The paired differences (MLP rank minus heuristic) are \rhval{analysis/analysis/test_41_70/d_mlprank_heur_t41/mean:2} ($p=$ \rhval{analysis/analysis/test_41_70/p_mlprank_heur_t41/mean}) at 41--70 and \rhval{analysis/analysis/test_41_70/d_mlprank_heur_t21/mean:3} ($p=$ \rhval{analysis/analysis/test_41_70/p_mlprank_heur_t21/mean}) at 21--40; at 11--20 both are at the ceiling. Relative to the MLP with $-\log p$ cost, the rank cost adds \rhval{analysis/analysis/test_41_70/d_mlprank_mlp_t41/mean:2} at 41--70 ($p=$ \rhval{analysis/analysis/test_41_70/p_mlprank_mlp_t41/mean}). This is the only configuration that beats the heuristic, it was not among the registered hypotheses, and it was selected after seeing the test bins from a small family of four policy/cost combinations; the result should be read as a hypothesis for follow-up rather than a confirmed advantage. It indicates that a learned ranking, used with a cost that does not penalise depth, can improve on the heuristic's ordering when the budget binds, even though the MLP's defer rate in Table~\ref{tab:diag} is lower than the heuristic's.

\paragraph{Learned instead of sinusoidal positions.} Replacing the sinusoidal encoding by learned absolute positions (positions beyond the training length are untrained, indices are clipped at the maximum) leaves the transformer's solved rate essentially unchanged: \rhval{abl_pos/transformer-learned-pos/test_11_20/solved/mean:2}, \rhval{abl_pos/transformer-learned-pos/test_21_40/solved/mean:2}, \rhval{abl_pos/transformer-learned-pos/test_41_70/solved/mean:2} for the three bins versus the main rows in Table~\ref{tab:main} (paired differences \rhval{analysis/analysis/test_41_70/d_lp_tf_t11/mean:3}, \rhval{analysis/analysis/test_41_70/d_lp_tf_t21/mean:3}, \rhval{analysis/analysis/test_41_70/d_lp_tf_t41/mean:3}; $p=$ \rhval{analysis/analysis/test_41_70/p_lp_tf_t11/mean}, \rhval{analysis/analysis/test_41_70/p_lp_tf_t21/mean}, \rhval{analysis/analysis/test_41_70/p_lp_tf_t41/mean}). Positional extrapolation is thus not the bottleneck we can detect; this does not exclude other forms of length-related failure in the encoder.

\paragraph{Training-set size and budget (Table~\ref{tab:sweeps}, Figs.~\ref{fig:ntrain}--\ref{fig:budget}).} Training-set size was varied with a fixed number of optimizer steps, so smaller sets are trained for more epochs. Contrary to the usual expectation, \emph{less} data extrapolated better: the transformer's solved rate at 11--20 falls from \rhval{analysis/analysis/test_41_70/nt1000_t11_m/mean:2} (1000 sequents) to \rhval{analysis/analysis/test_41_70/nt20000_t11_m/mean:2} (20000), and at 41--70 from \rhval{analysis/analysis/test_41_70/nt1000_t41_m/mean:2} to \rhval{analysis/analysis/test_41_70/nt20000_t41_m/mean:2}, with the largest step between 1000 and 4000 sequents (Table~\ref{tab:sweeps}). The 1000-sequent policy still stays below the heuristic in both bins. This sweep is exploratory (we did not register it, ran it only with the same steps and learning rate, and did not log in-distribution accuracy for the smaller sets), so we do not know whether the effect is due to overfitting to the larger sets' structure, to the number of epochs, to the particular states in the smaller sets, or to noise in the optimiser; it does show that the in-distribution data volume is not the limiting factor for size extrapolation here. For the budget sweep at 41--70, all three systems improve with budget; the heuristic's solved rate is \rhval{analysis/analysis/test_41_70/bud25_h_m/mean:2}, \rhval{analysis/analysis/test_41_70/bud50_h_m/mean:2}, \rhval{analysis/analysis/test_41_70/bud100_h_m/mean:2}, \rhval{analysis/analysis/test_41_70/bud200_h_m/mean:2} for budgets 25, 50, 100, 200, the transformer's \rhval{analysis/analysis/test_41_70/bud25_tf_m/mean:2}, \rhval{analysis/analysis/test_41_70/bud50_tf_m/mean:2}, \rhval{analysis/analysis/test_41_70/bud100_tf_m/mean:2}, \rhval{analysis/analysis/test_41_70/bud200_tf_m/mean:2}, and BFS's \rhval{analysis/analysis/test_41_70/bud25_b_m/mean:2}, \rhval{analysis/analysis/test_41_70/bud50_b_m/mean:2}, \rhval{analysis/analysis/test_41_70/bud100_b_m/mean:2}, \rhval{analysis/analysis/test_41_70/bud200_b_m/mean:2}. The ordering heuristic $>$ transformer $>$ BFS in mean solved rate holds at every budget tested; the transformer's gap to the heuristic does not close with larger budgets in this range.

\begin{table}[t]\centering\caption{Sweeps (solved rate, mean $\pm$ std over \rhval{const/n_seeds} seeds). Top: training sequents for the transformer (the 20000 column is the main-table row). Bottom: expansion budget at 41--70 connectives (the budget-100 column is the main-table row).}\label{tab:sweeps}
\resizebox{\columnwidth}{!}{\begin{tabular}{c}
\begin{tabular}{lcccc}\toprule
Training sequents & 1000 & 4000 & 10000 & 20000\\\midrule
bin 11--20 & \rhval{analysis/analysis/test_41_70/nt1000_t11_m/mean:2}$\pm$\rhval{analysis/analysis/test_41_70/nt1000_t11_sd/mean:2} & \rhval{analysis/analysis/test_41_70/nt4000_t11_m/mean:2}$\pm$\rhval{analysis/analysis/test_41_70/nt4000_t11_sd/mean:2} & \rhval{analysis/analysis/test_41_70/nt10000_t11_m/mean:2}$\pm$\rhval{analysis/analysis/test_41_70/nt10000_t11_sd/mean:2} & \rhval{analysis/analysis/test_41_70/nt20000_t11_m/mean:2}$\pm$\rhval{analysis/analysis/test_41_70/nt20000_t11_sd/mean:2}\\
bin 41--70 & \rhval{analysis/analysis/test_41_70/nt1000_t41_m/mean:2}$\pm$\rhval{analysis/analysis/test_41_70/nt1000_t41_sd/mean:2} & \rhval{analysis/analysis/test_41_70/nt4000_t41_m/mean:2}$\pm$\rhval{analysis/analysis/test_41_70/nt4000_t41_sd/mean:2} & \rhval{analysis/analysis/test_41_70/nt10000_t41_m/mean:2}$\pm$\rhval{analysis/analysis/test_41_70/nt10000_t41_sd/mean:2} & \rhval{analysis/analysis/test_41_70/nt20000_t41_m/mean:2}$\pm$\rhval{analysis/analysis/test_41_70/nt20000_t41_sd/mean:2}\\
\bottomrule\end{tabular}\\[4pt]
\begin{tabular}{lcccc}\toprule
Budget (bin 41--70) & 25 & 50 & 100 & 200\\\midrule
Transformer & \rhval{analysis/analysis/test_41_70/bud25_tf_m/mean:2}$\pm$\rhval{analysis/analysis/test_41_70/bud25_tf_sd/mean:2} & \rhval{analysis/analysis/test_41_70/bud50_tf_m/mean:2}$\pm$\rhval{analysis/analysis/test_41_70/bud50_tf_sd/mean:2} & \rhval{analysis/analysis/test_41_70/bud100_tf_m/mean:2}$\pm$\rhval{analysis/analysis/test_41_70/bud100_tf_sd/mean:2} & \rhval{analysis/analysis/test_41_70/bud200_tf_m/mean:2}$\pm$\rhval{analysis/analysis/test_41_70/bud200_tf_sd/mean:2}\\
Heuristic & \rhval{analysis/analysis/test_41_70/bud25_h_m/mean:2}$\pm$\rhval{analysis/analysis/test_41_70/bud25_h_sd/mean:2} & \rhval{analysis/analysis/test_41_70/bud50_h_m/mean:2}$\pm$\rhval{analysis/analysis/test_41_70/bud50_h_sd/mean:2} & \rhval{analysis/analysis/test_41_70/bud100_h_m/mean:2}$\pm$\rhval{analysis/analysis/test_41_70/bud100_h_sd/mean:2} & \rhval{analysis/analysis/test_41_70/bud200_h_m/mean:2}$\pm$\rhval{analysis/analysis/test_41_70/bud200_h_sd/mean:2}\\
BFS & \rhval{analysis/analysis/test_41_70/bud25_b_m/mean:2}$\pm$\rhval{analysis/analysis/test_41_70/bud25_b_sd/mean:2} & \rhval{analysis/analysis/test_41_70/bud50_b_m/mean:2}$\pm$\rhval{analysis/analysis/test_41_70/bud50_b_sd/mean:2} & \rhval{analysis/analysis/test_41_70/bud100_b_m/mean:2}$\pm$\rhval{analysis/analysis/test_41_70/bud100_b_sd/mean:2} & \rhval{analysis/analysis/test_41_70/bud200_b_m/mean:2}$\pm$\rhval{analysis/analysis/test_41_70/bud200_b_sd/mean:2}\\
\bottomrule\end{tabular}
\end{tabular}}

\end{table}
\begin{figure}[t]\centering\includegraphics[width=0.9\columnwidth]{sweep_ntrain_fig.pdf}\caption{Transformer solved rate vs.\ number of training sequents (mean $\pm$ std over seeds).}\label{fig:ntrain}\end{figure}
\begin{figure}[t]\centering\includegraphics[width=0.9\columnwidth]{sweep_budget_fig.pdf}\caption{Solved rate vs.\ expansion budget in the 41--70 bin (mean $\pm$ std over seeds).}\label{fig:budget}\end{figure}
